\documentclass[english]{article}

\usepackage{ae,aecompl}
\usepackage[T1]{fontenc}
\usepackage[latin9]{inputenc}
\usepackage{textcomp}
\usepackage{graphicx}
\usepackage{eurosym}
\usepackage{hyperref}
\usepackage{float}
\usepackage{fancyhdr}
\pagestyle{fancy}
\lhead{PowerEnJoy RASD}

\newcommand{\carsharing}{\textit {car sharing }}
\newcommand{\powerenjoy}{\textit{PowerEnJoy }}
\newcommand{\registereduser}{\textit {registered user }}
\newcommand{\staff}{\textit{staff }}
\newcommand{\service}{\textit{service }}
\newcommand{\safearea}{\textit{safe area }}
\newcommand{\powergrid}{\textit{Power Grid }}
\newcommand{\resevation}{\textit{reservation }}
\newcommand{\stand}{\textit{stand }}
\newcommand{\personalpin}{\textit{personal pin }}
\usepackage{tabto}
\usepackage{verbatim}

\makeatletter
\usepackage{babel}


\makeatother

\usepackage{babel}
\begin{document}
\begin{figure}
	\centering
	\includegraphics[scale=0.5]{logo.pdf} 
\end{figure}


\title{PowerEnJoy\\
 Requirement Analysis and Specification Document\\
}

\date{A.A 2016/2017}

\author{Erba Alessandro\\
 Leveni Filippo\\
 Lodi Luca}

\maketitle
\pagebreak{}

\tableofcontents{} \pagebreak{}

\section{Introduction}
	\subsection{Purpose }
	
		The purpose of our team is to project PowerEnjoy, which is  a platform that will be used to manage a car sharing service 
  in the city area.\\
  The platform will allow clients to find and reserve an electric car using mobile app from their self phone, and it will also    manage the communication of the emergency situations to the proper staff.\\
  Furthermore, a particular PowerEnjoy feature is the application of certain discounts to its users.\\
  \tab \tab This document aims to describe the high-level functionalities that will be offered by PowerEnjoy service. The    RASD is intended to be viewed by the stakeholders, that will evaluate the correctness of any assumption and decision
  in this document.
  
 \subsection{Scope and Problem Description}
 
  The government of a large city would like to introduce an ecological way to travel within the city area. In particular,     uniform distribution of cars in the city is a significant target.\\
  \tab \tab The cars are electric as anticipated and inside them it is available an integrated interface. Through this interface   is possible to communicate with the system (e.g. to call emergency staff).\\
  \tab \tab Users can choose and reserve a car through a mobile app. The system answers to the request by showing the    time left to the expiration of the reservation.\\
  The staff is divided in three categories:\\
   - Field staff\\
   - Emergency staff\\
   - Management staff\\
  \tab \tab Field staff use the mobile application to inform the system about their availability and to confirm that they are    going to take care of a certain call (e.g. retrieve a broken car, move a vehicle with a low battery or to ensure the uniform    distribution in the city, etc...).\\
  \tab \tab Emergency staff use the mobile application to answer customers calls and to take care of special situations (e.g.    no money in the costumer's credit card, car accident, etc...).\\
  \tab \tab Management staff use the mobile application to handle the service (e.g. set the safe areas, fares, discounts,     etc...).\\
  \tab \tab The system guarantees a uniform distribution of cars in the city. In particular, the city is divided in       circumferences and sectors. The system automatically computes the distribution of available cars both in the areas     between the circumferences and in the sectors, based on the GPS information it receives from each car. Uniformity is     guaranteed by ensuring the same amount of available cars for each area between the circumferences and for each sector.   \\
  \tab \tab A user can take advantage of the "Money saving" option selecting it on the integrated interface in the car. After    that he/she can input his/her final destination and the system provides information about the station where to leave the    car to get a discount. This destination is determined to ensure the uniform distribution of cars and depends both     on the destination of the user and on the availability of power plugs at the selected station.
	
	\subsection{Goals}
	We have selected a list of objectives that we think \powerenjoy  has to reach:
	\begin{description}
		
		\begin{comment} drivers for PowerEnjoy's success \end{comment}
		\item [G.1]{ To allow user to simply access to the \powerenjoy services. }
		\item [G.2]{ To guarantee an higher  probability of finding a car, compared to other car sharing services. }
		\item [G.3]{ To guarantee an efficient manutenution and recharge process of the cars, oriented to favorite their fairly distribution over the territory. }
		\begin{comment} user's goals \end{comment}
		\item [G.4]{ To allow users to make a reservation for a car in advance, and to cancel a reservation by its expiration time. }
		\item [G.5]{  To transparently charge the user for the service, minimizing its interaction with payment interfaces. }
		\item [G.6]{  To allow users to drive cars, minimizing their efforts for finding safe areas where to park.}
		\item [G.7]{ To encourage users' virtuous behaviour  in relation to service's fairness. }
		\begin{comment} service efficiency goals \end{comment}
		\item [G.8]{ To create a system that gives the possibility to be monitored and administrated by \powerenjoy's authorized personal. }
		\item [G.9]{  To empower the authorized personal in order to be able to quickly react to emergencies, potentially dangerous situations and users' requests for help.}
		\item [G.10]{ To create a system that provides a programmable interface, in order reduce the costs of further enanchements / manutenution to the system exploiting internal personal or third party developers.}
	\end{description}

	\subsection{Domain Assumptions }
	We suppose that the following conditions hold in the analysed world:
	\begin{description}
		
		\begin{comment} drivers for PowerEnjoy's success \end{comment}
		\item [D.1]{ The geographical area covered by the service is included in the coverage area of most common mobile communication technologies (3g, 4g) offered by main telecommunications companies. }
		\item [D.2]{ Only correctly registered users can access the functionalities provided by the system. }
		\item [D.3]{ Cars are able to send diagnostic for the manutention process in addition to realtime telemetrics and status of the battery.  }
		\item [D.]{ Cars can be selectively remotely locked, unlocked, powered on and powered off by authorized personal. }
		\item [D.]{  }
		\item [D.]{  }
		\item [D.]{  }
		\item [D.]{  }
		\item [D.]{  }
		\item [D.]{  }	

	\end{description}

	\subsection{Glossary}
		\subsubsection{Car Sharing}
			\carsharing is a model of car rental where people rent cars for short periods of time, often by the hour.
		\subsubsection{PowerEnJoy}
			\powerenjoy is the \carsharing brand object of this document. 
		\subsubsection{Registered User}
			A \registereduser is a person subscribed to \powerenjoy which can access to the \carsharing through his smart-phone.
		\subsubsection{Staff}
			\staff is the set of people that are \registereduser but can perform special operations. They are divided in three different categories.
			\begin{itemize}
				\item {Field Staff}
				\begin{itemize}
				\item They own a passepartout for cars that has all kind of issues, for example they can menage cars with low battery.
				\end{itemize}
				\item{Emergency Staff}
				\begin{itemize}
				\item They manage users issues. They are available to communicate with users.
				\end{itemize}
				\item{Management Staff}
				\begin{itemize}
				\item They manage the system, for example they can change system parameters. They can add new \safearea, modify \carsharing prices.
				\end{itemize}
			\end{itemize}
		\subsubsection{Service}
			The \service is the group of actions that a \registereduser can fulfill trough the mobile.	
		\subsubsection {Safe Area}
			A \safearea is a  geographical place\footnote{defined by a set of GPS positions} on the map in which parking is allowed. \safearea are saved into the system. A \registereduser can end a \carsharing parking his car only in safe areas.
	\subsubsection{Power Grid}
		A \powergrid is the station that charge the battery of an electric car. All \powergrid are displayed in the map on-board \powerenjoy cars. Each \powergrid is in a \safearea. A \registereduser who plugs his car into a \powergrid receives the discount expected within \powerenjoy rules.
	\subsubsection{Reservation}
		A \resevation is the possibility that a \registereduser has to book a car at the latest for one hour.
	\subsubsection{Stand}
		A \stand is the possibility that a \registereduser has to pause his \carsharing. During a \stand the \registereduser can leave the car, than he can re-unlock it through his smart-phone.
	\subsubsection{Personal  Personal Identification Number (pin)}
		 \personalpin is a number of 5 digits that the user inserts on the car screen to enable driving commands.

\subsection{Further Developments}
	We think that one of the most important things is to release a Service that can be extended and improved. It's trivial that a service like this must be able to be change and be updated.
	For this purpose we will develop a system open to new extensions. \\
	Here we provide a series of possible extensions that could be necessary in the future.
	E.g.:
	\begin{itemize}
		\item New kind of discount policies can be introduced to the system, for example discounts for new users. 
		\item New system policies may be necessary, for example new \carsharing conditions.
		\item New kind of vehicles can be introduced, for example electric scooter and electric bikes.
		\item Support to new payment methods, for example Paypal\textregistered, NFC\footnote{Near Field Communication; NFC devices are used in contactless payment systems, similar to those used in credit cards and electronic ticket smartcards and allow mobile payment to replace/supplement these systems. \href{https://en.wikipedia.org/wiki/Near_field_communication}{NFC on Wikipedia}}.
		\item Develop new services in order to increase the user experience, for example it may be useful to allow communication among users in order to share a car ride and save money.
	\end{itemize}
	\subsection{Used Tools}
	\begin{itemize}
		\item \LaTeX\\
		\item GitHub\\
	\end{itemize}
\section{Specific Requirements}
	\subsection{Functional Requirements}
	\subsection{Non Functional Requirements}

\section{Scenarios Identifying}
\section{UML Models And Use Cases}
	\subsection{Use Cases Diagram}
	\subsection{Actors Identifying}
	\subsection{Use Cases}
	
	\subsubsection{Search Car}
	\begin{center}
	\begin{tabular}{l||p{10cm}}
	\textbf{Name} 
		& Search Car\\ [8px]
	\textbf{Actors} 
		& PowerEnJoy User\\ [8px]
	\textbf{Assumptions} 
		& \begin{itemize}
			\item The user is logged in to the system.
		\end{itemize}\\
	\textbf{Flow of events}
		& \begin{enumerate}
	 		\item The user opens PowerEnJoy app on his smartphone 
	 		\item The user interface displays all available cars around the user position according to the setted range. 
			\item Optionally the user can search the cars around an inserted address, or near a remarkable place. 
		\end{enumerate}\\ 
	\textbf{Exit conditions}
		&The user finds a car that satisfies his need and choses too reserve it.\\ [8px]
	\textbf{Exceptions}
		& \begin{itemize}
			\item The user does not find an useful car and exit the application.
			\item The user want not really reserve a car. 
		\end{itemize}
	\end{tabular}
	\end{center}

	\subsubsection{Reserve Car}
	\begin{center}
	\begin{tabular}{l||p{10cm}}
	\textbf{Name} 
		& Reserve Car\\ [8px]
	\textbf{Actors} 
		& PowerEnJoy User\\ [8px]
	\textbf{Assumptions} 
		& \begin{itemize}
			\item The user is logged in to the system.
			\item The user has Searched Cars and has find an optimal one
		\end{itemize}\\
	\textbf{Flow of events}
		& \begin{enumerate}
 			\item The user selects the car that he want reserve
 			\item The user interface displays all informations about the car (E.g. Autonomy in kilometers) and a button for reserve it. 
			\item The user choses to reserve the car. 
			\item The user confirms his choose.
		\end{enumerate}\\ 
	\textbf{Exit conditions}
		&The system confirms the successful reservation, displays the time before reservation expiration and the route for reach it from the current position of the user. \\ [8px]
	\textbf{Exceptions}
		& \begin{itemize}
			\item The user exits the procedure before confirming it.
			\item Another user reserves another car while our user is reserving it, the system notify the user. The reservation ends unsuccessfully.
		\end{itemize}
	\end{tabular}
	\end{center}

	\subsubsection{Cancel Reservation}
	\begin{center}
	\begin{tabular}{l||p{10cm}}
	\textbf{Name} 
		& Cancel Resarvation\\ [8px]
	\textbf{Actors} 
		& PowerEnJoy User\\ [8px]
	\textbf{Assumptions} 
		& \begin{itemize}
			\item The user is logged in to the system.
			\item The user has an active reservation.
		\end{itemize}\\
	\textbf{Flow of events}
		& \begin{enumerate}
 			\item The user opens PowerEnJoy app on his smartphone 
 			\item The user interface displays user's current reservation.
			\item The user eliminates the reservation clicking on the cancel reservation button.
			\item The user confirms his selection.
		\end{enumerate}\\ 
	\textbf{Exit conditions}
		&The user successfully cancel his reservation. The reserved car returns available for other users.\\ [8px]
	\textbf{Exceptions}
		&...\\[8px]
	\end{tabular}
	\end{center}

	\subsubsection{Pick up car}
	\begin{center}
	\begin{tabular}{l||p{10cm}}
	\textbf{Name} 
		& Pick up car\\ [8px]
	\textbf{Actors} 
		& PowerEnJoy User\\ [8px]
	\textbf{Assumptions} 
	& \begin{itemize}
		\item The user is logged in to the system.
		\item The user has a reservation active.
		\item The user is no far than 10 meters from the car
		\item The user inserts the correct \personalpin 
	\end{itemize}\\
	\textbf{Flow of events}
		& \begin{enumerate}
 		\item The user is close to his reserved car.
 		\item The user selects the unlock car button (displayed in the application) UI on it's smartphone.
		\item The system unlocks the car.
		\item The user gets in the car and insets his \personalpin.
		\item The system detects all passengers from cars sensors.
		\item Optionally the user can decide to activate the "Money Saving Option"
		\item The user starts driving his car.
		\item the system calculates the fare and shows it on the display
		\end{enumerate}\\ 
	\textbf{Exit conditions}
		&\begin{itemize}
			\item The user makes a \stand.
			\item The user releases car.
		\end{itemize}\\
	\textbf{Exceptions}
		& \begin{itemize}
			\item The user run out of money on his credit card.
			\item The car runs out of battery.
		\end{itemize}
	\end{tabular}
	\end{center}
	
	\subsubsection{Stand}
	\begin{center}
	\begin{tabular}{l||p{10cm}}
	\textbf{Name} 
		& Stand\\ [8px]
	\textbf{Actors} 
		& PowerEnJoy User\\ [8px]
	\textbf{Assumptions} 
	& \begin{itemize}
		\item The user has picked up a car. 
	\end{itemize}\\
	\textbf{Flow of events}
		& \begin{enumerate}
 		\item The user parks the car.
 		\item The user selects the\stand option on the car screen.
		\item Driving command are disabled by the system. 
		\item The user get off the car.
		\item The system locks the car.
		\end{enumerate}\\ 
	\textbf{Exit conditions}
		&\begin{itemize}
			\item When the user wants end the \stand he re-unlocks and start driving the car as he has done in the "Pick up car" use case.
		\end{itemize}\\
	\textbf{Exceptions}
		& \begin{itemize}
			\item The user run out of money on his credit card. (The user will not be able to re-unlock the car.)
		\end{itemize}
	\end{tabular}
	\end{center}
	
	\subsubsection{Money Saving Option}
	\begin{center}
	\begin{tabular}{l||p{10cm}}
	\textbf{Name} 
		& Money Saving Option\\ [8px]
	\textbf{Actors} 
		& PowerEnJoy User\\ [8px]
	\textbf{Assumptions} 
	& \begin{itemize}
		\item The user has picked up a car. 
	\end{itemize}\\
	\textbf{Flow of events}
		& \begin{enumerate}
 		\item The user selects his destination address on the car screen
		\item The system purposes the \powergrid station according to algorithm constraints.
		\end{enumerate}\\ 
	\textbf{Exit conditions}
		&\begin{itemize}
			\item The user releases the car.
		\end{itemize}\\
	\textbf{Exceptions}
		& \begin{itemize}
			\item There are no available \powergrid stations in 1 km range from the destination. (The system during driving will notify the user through the car screen if a \powergrid becomes free.)
		\end{itemize}
	\end{tabular}
	\end{center}

	\subsubsection{Release Car}
	\begin{center}
	\begin{tabular}{l||p{10cm}}
	\textbf{Name} 
		& Release Car\\ [8px]
	\textbf{Actors} 
		& PowerEnJoy User\\ [8px]
	\textbf{Assumptions} 
	& \begin{itemize}
		\item The user has picked up a car. 
	\end{itemize}\\
	\textbf{Flow of events}
		& \begin{enumerate}
 		\item The user parks the car in a \safearea.
		\item The user selects the end \carsharing option on the car screen.
		\item Driving command are disabled by the system. 
		\item The system starts calculating fare discounts.
		\end{enumerate}\\ 
	\textbf{Exit conditions}
		&\begin{itemize}
			\item The user get off the car, than the system locks the car.
		\end{itemize}\\
	\textbf{Exceptions}
		& \begin{itemize}
			\item The user is not in a safe area and he cannot end \carsharing. 
		\end{itemize}
	\end{tabular}
	\end{center}

\subsubsection{Plug car}
	\begin{center}
	\begin{tabular}{l||p{10cm}}
	\textbf{Name} 
		& Release Car\\ [8px]
	\textbf{Actors} 
		& PowerEnJoy User\\ [8px]
	\textbf{Assumptions} 
	& \begin{itemize}
		\item The user has parked a car near a \powergrid and has ended his \carsharing. 
	\end{itemize}\\
	\textbf{Flow of events}
		& \begin{enumerate}
 			\item The user inserts his \personalpin into the \powergrid screen, in order to let the \powergrid enabled to charge cars.
			\item The user plugs the car to the \powergrid.
			\item The system detects that the car of our user is charging at the \powergrid unlocked by the same user. 
		\end{enumerate}\\ 
	\textbf{Exit conditions}
		&\begin{itemize}
			\item The user receives the expected discount.
		\end{itemize}\\
	\textbf{Exceptions}
		&...\\[8px]
	\end{tabular}
	\end{center}

	\subsection{Class Diagram}
\section{External Interfaces}
\section{Alloy Model}
\section{Hours of Work}

\end{document}